---
title: Termination Under Bounded Delegation
category: algorithms
slug: termination-bounded-delegation
nav_order: 3
status: current
---

\section{Termination Under Bounded Delegation}

\subsection*{Implementation Constants}

The proof below is grounded in two concrete bounds present in the codebase:

\begin{itemize}
  \item \textbf{Supervisor turn bound}: $D_S = 200$ — \texttt{max\_turns} in
        \texttt{process\_conversation} (\texttt{orchestrator\_core.py})
  \item \textbf{Worker execution timeout}: $T_W = 1200\text{s}$ —
        \texttt{\_WORKER\_RUN\_TIMEOUT} in \texttt{DelegationMixin}
\end{itemize}

Both bounds are finite. Together they guarantee the system cannot run indefinitely.

\subsection*{Definitions}

Let $D_S \in \mathbb{N}$ be the maximum supervisor turn count.

Let $T_W \in \mathbb{R}^+$ be the maximum wall-clock time per worker run.

Let the supervisor state machine include:

\[ \{Delegating,\ Awaiting,\ Evaluating,\ Synthesizing,\ Terminated\} \]

Let:

\[ f: C \rightarrow \{complete,\ incomplete\} \]

be the evaluation function — in the implementation, this is the LLM's decision
to emit a tool call (incomplete) or a terminal text reply (complete).

\subsection*{Transition Rules}

\[ Evaluating \rightarrow Delegating \quad \text{iff } f(C) = incomplete \text{ and } d < D_S \]

\[ Evaluating \rightarrow Synthesizing \quad \text{iff } f(C) = complete \text{ or } d \geq D_S \]

Each turn increments the counter:

\[ d := d + 1 \]

\subsection*{Bounded Constraint}

\[ d \leq D_S = 200 \quad \text{and} \quad T_{worker} \leq T_W = 1200\text{s} \]

\subsection*{Theorem}

Given $D_S < \infty$ and $T_W < \infty$, total execution terminates.

\subsection*{Proof}

\textbf{Worker termination:} Each worker run either completes or is interrupted by
\texttt{\_wait\_for\_run\_completion} after $T_W = 1200\text{s}$.
\texttt{asyncio.TimeoutError} is caught, \texttt{execution\_had\_error} is set,
and the delegation loop continues. Workers cannot run indefinitely.

\textbf{Supervisor termination:} Each iteration of the \texttt{while turn\_count < max\_turns}
loop in \texttt{process\_conversation} strictly increments \texttt{turn\_count}.
After at most $D_S = 200$ iterations the loop exits.

\textbf{Composition:} Since each worker terminates within $T_W$ and the supervisor
executes at most $D_S$ worker delegations, total execution terminates within:

\[ T_{total} \leq D_S \cdot T_W = 200 \times 1200\text{s} = 240{,}000\text{s} \]

No infinite execution path exists.

\[ \therefore \text{Termination is guaranteed within } D_S \cdot T_W \text{ seconds.} \]

\subsection*{Known Failure Mode}

When $d = D_S$, \texttt{process\_conversation} yields:

\begin{verbatim}
{"type": "error", "content": "Maximum conversation turns exceeded"}
\end{verbatim}

and stamps \texttt{failed\_at}. There is no re-planning, no escalation,
and no partial synthesis. The system terminates without success and does not
signal recovery. This is the primary open gap. See \textit{Liveness Gap Analysis}.

\qed